\documentclass[11pt]{article}

\usepackage{graphicx}
\usepackage{amsmath}
\usepackage[margin=1in]{geometry}
\usepackage{parskip}

\title{Time Perception from Interacting Oscillators\\\large An Overnight Report}
\author{Exploratory session}
\date{2026-06-11}

\begin{document}
\maketitle

\section*{What this project is}

The goal is a small, self-contained model of \emph{subjective} time --- the feeling
of how much time has passed --- built entirely out of \textbf{oscillators} and tested
as a working program. No language model is involved.

An \textbf{oscillator} here is just something that cycles steadily, like the second
hand of a clock going round and round. Two numbers describe it at any moment: its
\textbf{frequency} (how fast it goes round --- cycles per second) and its
\textbf{phase} (where in the current cycle it is right now, which we can think of as
an angle from $0^\circ$ to $360^\circ$, or equivalently a fraction from $0$ to $1$).
A collection of such oscillators is a ``bank.''

The central idea is that a bank of oscillators running at different speeds can act as
a clock, and that the \emph{feeling} of time can be read off that clock in different
ways depending on what you ask. This report describes what was built on top of the
Stage~1 engine, what the experiments showed (including two places where the first
guess was wrong and had to be corrected), and where to go next.

\section{The six main results, in one paragraph each}

\textbf{1. The ``holiday paradox'' works, and it is not fragile.}
The same system makes an idle stretch of time feel \emph{long} while you are living
through it, yet makes a busy stretch be \emph{remembered} as having lasted longer.
These two facts point in opposite directions, and both fall out of the model.
Crucially, this is not a fluke of one lucky setting --- it holds across a wide range
of conditions (Section~\ref{sec:dissoc}).

\textbf{2. The scale-invariance of timing errors needs a specific kind of noise.}
Human time estimates get \emph{proportionally} more variable for longer intervals
(your error timing 10 seconds is about ten times bigger than your error timing
1 second). This appears only when the random fluctuation in the clock is
\emph{shared} across all the oscillators at once, not when each oscillator wobbles on
its own. The kind of randomness matters, not just the amount
(Section~\ref{sec:laws}).

\textbf{3. The textbook ``psychophysics'' of timing does not come from the
oscillators themselves --- it needs an extra step.} Two classic laboratory findings
about time perception turned out to require representing time on a \emph{ratio}
(logarithmic) scale, which the raw oscillator machinery does not do by itself.
Oscillations are necessary but not sufficient (Section~\ref{sec:laws}).

\textbf{4. Making the oscillators influence each other helps only a little, and too
much of it is catastrophic.} When the oscillators pull on one another, a small amount
improved one useful signal, but past a tipping point they all fell into lockstep and
the clock stopped being able to tell different times apart (Section~\ref{sec:couple}).

\textbf{5. There are two ways to build a multi-speed clock, with opposite
weaknesses.} One is compact and exact but shatters under the slightest noise; the
other is bulky and approximate but degrades gracefully. This is a genuine engineering
trade-off between \emph{capacity} and \emph{robustness} (Section~\ref{sec:two}).

\textbf{6. That trade-off can be fixed by being selective about which oscillators talk
to each other.} Letting each oscillator influence only its immediate neighbour (rather
than all of them at once) keeps the compact clock from shattering, without causing the
lockstep collapse from result~4. \emph{Who} is coupled to \emph{whom} matters more
than how strongly (Section~\ref{sec:struct}).

\section{How the model is put together}

There is one engine and several different ``read-outs'' --- different ways of
interrogating the same underlying state.

\begin{verbatim}
  incoming events
        |
        v
  [ attention gate ]      <- decides how much each event "counts"
        |
        v
  bank of oscillators  -----------------+
        |                               |
        v                               v
  three running tallies            the pattern of phases
  --------------------             ---------------------
  - felt time (now)                used to estimate "how
  - remembered amount              long was that?", and
  - recency                        to compare durations
\end{verbatim}

The pieces, in plain terms:

\begin{itemize}
\item \textbf{Event:} something happening --- a stimulus arriving, a message, an
action. In the model an event gives the oscillators a little nudge.
\item \textbf{Attention gate:} a control that decides how strongly incoming events
register. When a lot is going on, attention is absorbed by the events and the gate to
the internal clock partly closes; when little is going on, the gate is open and the
clock is watched closely. (This is the standard explanation for why a watched pot
never boils --- when you attend \emph{to time itself}, it feels slow.)
\item \textbf{Felt time (``now''):} a running tally of subjective time as it is being
lived.
\item \textbf{Remembered amount:} a running tally of how much happened, used when
looking \emph{back} on a stretch of time.
\item \textbf{Recency:} a signal that fades as time passes since the last event,
letting the system sense ``how long since something happened.''
\end{itemize}

\section{The holiday paradox}
\label{sec:dissoc}

The headline result. Take two stretches of clock time, each 60 seconds long. In one
(``busy'') many events happen; in the other (``idle'') almost nothing happens. The
model reports:

\begin{center}
\begin{tabular}{llll}
\hline
What we measure & Busy stretch & Idle stretch & Meaning \\
\hline
Felt time, as lived & $\sim$6 seconds & $\sim$43 seconds & the idle stretch \textbf{drags} \\
Remembered amount & high ($\sim$4300) & low ($\sim$340) & the busy stretch is recalled as \textbf{longer} \\
\hline
\end{tabular}
\end{center}

So the busy stretch \emph{flies by} in the moment but is \emph{remembered as long};
the idle stretch \emph{drags} in the moment but \emph{shrinks to almost nothing} in
memory. This is the everyday ``holiday paradox'': an absorbing holiday races past
while you are on it, yet feels long when you look back, because so much happened. The
model reproduces both halves because the two read-outs are driven by opposite things
--- felt time runs fast when little is happening (attention is free to watch the
clock), while the remembered amount simply counts up how much occurred. This is the
reason the system is more than a stopwatch: a stopwatch would give the same answer
both ways.

\begin{figure}[ht]
\centering
\includegraphics[width=\textwidth]{stage1_results.png}
\caption{Left: felt (prospective) time vs.\ real time --- the idle stretch tracks
near real time (drags), the busy stretch lags far behind (flies). Middle: a
``check-in'' fires when \emph{felt} idle time crosses a threshold. Right: the
scale-invariance of error (see Section~\ref{sec:laws}).}
\end{figure}

\subsection{Is it just a lucky setting?}

A fair worry: maybe the paradox only appears because of two convenient cases. To
check, the event rate was varied smoothly from almost-nothing to very-busy.

\begin{itemize}
\item Felt time fell \emph{smoothly and continuously} as things got busier (from
$\sim$55 down to $\sim$3 seconds).
\item Remembered amount rose \emph{smoothly and continuously} over the same range.
\end{itemize}

The two curves cross and keep pointing in opposite directions across the whole range,
and this stayed true for every setting of the gate's two tuning knobs that was tried.
In other words, the effect is a broad, stable regime, not a coincidence that holds
only at one exact balance point. (A ``knife-edge'' result is one that only works at a
single precise tuning and vanishes if you nudge it; this is not one.)

\begin{figure}[ht]
\centering
\includegraphics[width=\textwidth]{robustness_results.png}
\caption{Left: as events get more frequent, felt time (blue) falls while remembered
amount (red) rises --- the dissociation holds across the whole range. Middle and
right: ``busy flies'' holds for every setting of the two gate parameters; they only
change the steepness.}
\end{figure}

\section{Two classic timing laws --- and where they actually come from}
\label{sec:laws}

Psychologists have catalogued several robust facts about how people judge durations.
Three were tested directly against the oscillator clock. One came out naturally (with
a caveat); two did \textbf{not}, and understanding why was the most useful part of the
work.

Two definitions used throughout:

\begin{itemize}
\item \textbf{Coefficient of variation (CV):} the spread of a set of measurements
divided by their average. It is \emph{relative} error. A CV of $0.1$ means the typical
error is 10\% of the thing being measured, whatever its size.
\item \textbf{Noise / jitter:} small random variation added to the model. Two kinds.
\emph{Shared} (common-mode) jitter means all oscillators speed up or slow down
together on a given trial --- like the whole clock running slightly fast that day.
\emph{Independent} jitter means each oscillator wobbles on its own, unrelated to the
others.
\end{itemize}

\subsection{Scale-invariance of error (``Weber's law'')}

The fact: a person's error in judging a duration grows in \emph{proportion} to the
duration, so the \emph{relative} error (the CV) stays roughly constant across
durations. This proportional-error property is called the \textbf{scalar property} or
\textbf{Weber's law}.

\begin{itemize}
\item With \textbf{shared} jitter, the CV stayed flat at about $0.05$ across durations
from 2 to 32 seconds. That is exactly Weber's law, and it matches the textbook account
in which the clock's \emph{overall rate} varies a little from trial to trial.
\item With \textbf{independent} jitter, the CV did \emph{not} stay flat --- it grew
from $0.03$ to $0.23$ as durations got longer. Each oscillator drifting on its own
scrambles the clock's pattern at long durations rather than simply rescaling it.
\end{itemize}

The lesson: Weber's law pins down the \emph{kind} of randomness in the clock. It must
be a shared, whole-clock fluctuation, not independent per-oscillator wobble.

\subsection{Comparing two durations (``bisection'')}

The task: a person learns a ``short'' reference (say 4 seconds) and a ``long''
reference (16 seconds), then judges whether a new probe duration is closer to short or
long. The interesting finding is \emph{where the cross-over sits} --- the probe judged
``long'' half the time. In people, it lands near the \textbf{geometric mean} of the
two references.

The geometric mean of two numbers is the square root of their product:
$\sqrt{4\times16}=8$. It is the natural midpoint on a \emph{ratio} scale ($8$ is the
same factor ($\times2$) above $4$ as $16$ is above $8$), as opposed to the ordinary
\textbf{arithmetic mean}, $(4+16)/2=10$, which is the midpoint on an additive scale.

The raw oscillator comparison \textbf{fails} at this task. With only two stored
references, a probe sitting between them is too far from \emph{both} to match either,
and the decision collapses into noise (the cross-over landed at a meaningless
$\sim$4.7 seconds). The reason is that the oscillator bank's ability to distinguish
two times is set by a \emph{fixed absolute resolution} (governed by its fastest
oscillator) --- it has no built-in sense of \emph{ratios}, which is what the geometric
mean is about.

The cross-over only lands correctly (at $8.0$) once a read-out is added that
represents time on a \textbf{logarithmic scale} --- i.e.\ by the \emph{logarithm} of
the duration rather than the duration itself. On a log scale, equal \emph{ratios}
become equal \emph{distances}, so the ordinary midpoint of the log values is the
geometric mean of the durations. This is the extra step the oscillators do not provide
on their own.

\subsection{Estimates pulled toward the middle (``Vierordt's law'')}

The fact: when judging a range of durations, people overestimate the short ones and
underestimate the long ones, as if every estimate is pulled toward the middle of the
range. This central-pull is called \textbf{Vierordt's law}.

The standard explanation is \textbf{Bayesian}: a sensible observer combines a noisy
measurement with prior expectations about what durations are likely, and that prior
pulls estimates inward. Adding this Bayesian step on top of the log-scale read-out
reproduces the law cleanly. Plotting estimated against true duration gives a line
shallower than reality (a slope of $0.79$ instead of $1.0$; less than~1 means the
estimates are compressed toward the centre), with the cross-over from over- to
under-estimation sitting right at the middle of the range.

\begin{figure}[ht]
\centering
\includegraphics[width=\textwidth]{stage1b_results.png}
\caption{Top row, the raw oscillator clock: (A1) the CV is flat only with shared
noise; (A2) it cannot do the comparison task; (A3) its resolution is a fixed absolute
width, with no sense of ratio. Bottom row, the logarithmic read-out: (B1) flat CV,
(B2) cross-over at the geometric mean, (B3) central-pull estimation.}
\end{figure}

\paragraph{What these add up to.} The oscillator bank supplies the multi-speed clock
and, with shared jitter, the proportional error of Weber's law. But the
\emph{ratio-based} phenomena --- geometric-mean comparison and central-pull estimation
--- live in a logarithmic read-out layered on top, not in the oscillators themselves.

\section{Letting the oscillators influence each other}
\label{sec:couple}

So far the oscillators run independently. The natural next question is what happens
when they \emph{interact}, each pulling on the others. The standard textbook model was
used (\textbf{Kuramoto coupling}: every oscillator is nudged toward the average phase
of the group, with a strength called $K$). Two new terms:

\begin{itemize}
\item \textbf{Coherence (order parameter):} a single number from $0$ to $1$ measuring
how aligned the oscillators are. $0$ means scattered all around the cycle; $1$ means
all bunched at the same phase, moving as one.
\item \textbf{Synchronization:} what happens at strong coupling --- the oscillators
stop running at their own speeds and fall into lockstep at a common speed.
\end{itemize}

Three things were measured as the coupling strength $K$ increased:

\begin{itemize}
\item \textbf{The synchronization tipping point.} Up to about $K\approx8$ the
oscillators stayed mostly independent (low coherence); above it they snapped into
lockstep (coherence jumped toward $1$).
\item \textbf{How long the clock stays useful (its ``coding horizon'').} After an
event resets the oscillators to a common start, the pattern of phases is unique for a
while, then eventually repeats --- and once it repeats, the clock can no longer tell
those two times apart. The time until that first repeat is the \textbf{coding horizon}
(longer is better). Below the tipping point it was long (beyond the 30-second test
window). At the tipping point it \emph{collapsed} --- from over 30 seconds to under
half a second --- because once everything is in lockstep at one speed, the pattern
just repeats every cycle.
\item \textbf{The recency signal.} Right after an event the oscillators are aligned
(coherence near $1$), and coherence then fades as they drift apart --- that fading
\emph{is} a ``time since last event'' signal. Without coupling it faded very fast
(coherence only $0.10$ two seconds later). A little coupling made it fade more slowly.
But strong coupling froze it near $1$ forever, which is useless --- a signal that
never changes tells you nothing.
\end{itemize}

\paragraph{Verdict.} Coupling earns its keep only in a narrow window of \emph{weak}
coupling (around $K\approx6$--$8$), where it stretches the recency signal to last
about three times longer while the clock is still intact. Push past the tipping point
and the whole multi-speed clock collapses into a single useless rhythm.

\begin{figure}[ht]
\centering
\includegraphics[width=\textwidth]{stage2_results.png}
\caption{Left: the recency signal (coherence after a reset) for several coupling
strengths --- too little fades instantly, too much never fades. Middle: the
synchronization tipping point near $K\approx8$. Right: the coding horizon collapses at
the same point.}
\end{figure}

\section{Two ways to build a multi-speed clock}
\label{sec:two}

There are two quite different ways to get a clock that works across many timescales
from a handful of oscillators, with opposite strengths and weaknesses.

\textbf{The cascade (a stack of gears).} Stack oscillators like the gears of a
mechanical clock or the wheels of a car's \textbf{odometer}: each one runs exactly 10
times slower than the one below it. The fast wheel gives the fine detail; each slower
wheel adds another digit. This is wonderfully compact: just 4 such wheels cover a
1000-second range \emph{exactly}, the way 4 odometer wheels count to 9999. A linear
number of oscillators buys an exponential range.

\textbf{The incommensurate bank (many unrelated speeds).} Alternatively, use a
\emph{bank} of oscillators whose frequencies are \textbf{incommensurate} --- they
share no common beat, e.g.\ $1.0, 1.7, 2.9$ cycles per second rather than the neat
$1, 10, 100$ of the gears. No single one repeats in step with the others for a long
time, so their combined pattern stays unique over a long span.

Both were tested against noise (random error added to each oscillator's phase):

\begin{itemize}
\item The \textbf{cascade} is \emph{exact but brittle}. With no noise it is perfect.
But the tiniest disturbance is catastrophic: like a slipped odometer wheel, a small
error in a slow (``high-digit'') oscillator throws the reading off by a huge amount.
Just 1\% phase noise already produced 1-second errors; 2\% produced 10-second errors.
\item The \textbf{incommensurate bank} is \emph{approximate but graceful}. Over the
same range and the same noise, its error stayed tiny throughout (a few hundredths of a
second). The cost is that it needs more oscillators and covers a shorter unambiguous
range.
\end{itemize}

So: \textbf{cascade = compact, exact, fragile; bank = bulky, approximate, sturdy.}
This is a real capacity-versus-robustness trade-off, and it is the same lesson as the
coupling result --- the arrangement that packs in the most capability is also the most
easily broken.

\begin{figure}[ht]
\centering
\includegraphics[width=\textwidth]{cascade_results.png}
\caption{Left: the cascade decodes a long range exactly, where a single oscillator
just repeats (aliases). Right: under increasing phase noise the cascade's error
explodes (digit slips) while the incommensurate bank stays tiny.}
\end{figure}

\section{Fixing the trade-off: be selective about coupling}
\label{sec:struct}

Section~\ref{sec:couple} showed that coupling \emph{all} the oscillators together
destroys the clock. Section~\ref{sec:two} showed the compact cascade is fragile. The
resolution: couple the oscillators, but only \textbf{locally} --- let each gear lock
onto just its immediate faster neighbour, keeping the exact 10:1 ratio between them,
rather than everyone pulling on everyone.

The fast oscillator is the most reliable (it completes many cycles, so small drifts
average out), and each slower one can be gently kept in step with it, correcting drift
before it causes an odometer-style digit slip.

Testing with a 3\% speed drift added to each stage over a 60-second run:

\begin{itemize}
\item \textbf{No coupling:} typical error $7.25$ seconds (the digit-slips above).
\item \textbf{Local neighbour-coupling:} typical error $0.08$ seconds --- about a
90-fold improvement.
\item \textbf{Even very strong local coupling:} still $0.08$ seconds --- it does
\emph{not} collapse the clock, unlike the all-to-all coupling of
Section~\ref{sec:couple}.
\item \textbf{Tolerance to drift:} with no coupling, error grew to $\sim$18 seconds at
12\% drift; with local coupling it stayed under a third of a second throughout.
\end{itemize}

This is the satisfying synthesis: the brittle-but-compact cascade becomes \emph{robust}
when stabilised by \emph{local} coupling. The thing that matters is not how strong the
coupling is but \emph{who is coupled to whom} --- neighbour-to-neighbour coupling
preserves the ladder of speeds, whereas all-to-all coupling erases it. (One honest
blemish: occasional $\sim$4-second glitches remain at the exact moments a digit rolls
over.)

\begin{figure}[ht]
\centering
\includegraphics[width=\textwidth]{structured_results.png}
\caption{Left: decode error over time --- unlocked drifts and jumps, local locking
stays flat. Middle: error drops sharply at modest coupling and, unlike global
coupling, does not collapse even when strong. Right: local locking tolerates far more
drift.}
\end{figure}

\section{The self-monitoring loop}
\label{sec:selfreg}

A small but pleasing extra. The system raises a ``check-in'' flag after enough
\emph{felt} idle time. Feeding each check-in back in as if it were an event closes the
loop and produces a self-sustaining rhythm --- in the language of dynamical systems, a
\textbf{limit cycle} (a stable repeating pattern the system settles into on its own):

\begin{itemize}
\item Left alone with nothing happening, the system checks in at very regular
intervals (every $9.8$ seconds, extremely steady).
\item The interval ($9.8$\,s) is a bit longer than the raw threshold ($8$\,s), because
each check-in briefly busies the system and partly closes the attention gate, so felt
time accrues slowly until things settle --- the period is an \emph{emergent}
compromise, not simply the threshold.
\item Background activity \emph{suppresses} the rhythm smoothly: the busier things are,
the less often the system checks in (down to never, once enough is going on).
\end{itemize}

This is a concrete, dynamic version of the ``nested feedback loops'' idea: a rhythm the
system generates itself and that its context modulates.

\begin{figure}[ht]
\centering
\includegraphics[width=\textwidth]{selfreg_results.png}
\caption{Left: check-ins (blue) against external events (red) at several background
rates --- regular when idle, suppressed when busy. Right: the self-check rate falls
smoothly as the background gets busier.}
\end{figure}

\section{Honest limitations}

\begin{itemize}
\item \textbf{Weber's law is partly built in, not discovered.} The proportional error
comes from a shared, whole-clock fluctuation put in by hand. It is the standard and
correct modelling choice, but an assumption, not something the dynamics produced on
their own.
\item \textbf{The ratio-based laws are not oscillatory.} Geometric-mean comparison and
central-pull estimation come from the logarithmic Bayesian read-out, not from the
oscillators. The pure ``everything is oscillations'' ambition hits a genuine wall here.
\item \textbf{The ``remembered amount'' is a crude stand-in} (a sum of how much each
event disturbed the clock). It gets the \emph{direction} of the holiday paradox right
but is not calibrated against real data.
\item \textbf{The arousal/attention effect is a knob, not an explanation.} Turning it
shifts estimates the right way, but that only confirms the dial is wired up sensibly.
\item \textbf{The coupling experiments use a moderate spread of speeds} so the tipping
point is reachable in simulation. The qualitative result does not depend on that
choice, but the specific numbers do.
\end{itemize}

\section{Where to go next (in priority order)}

\begin{enumerate}
\item \textbf{A fully self-correcting clock.} Combine the local-coupling stabiliser
with a \emph{noisy fastest oscillator} (right now the fast one is assumed perfect) and
add a little redundancy to remove the digit-rollover glitches.
\item \textbf{One unified pipeline.} Wire the oscillator clock, the logarithmic
read-out, and the Bayesian step into a single object, so one model produces all the
timing behaviours together.
\item \textbf{Couple the self-check rhythm to the ``remembered amount,''} so that a
\emph{full} stretch of time changes the rhythm, not just an empty one.
\item \textbf{Ground ``how much an event counts'' in real surprise} (how unexpected an
input is) if and when this is reconnected to a language model.
\end{enumerate}

\section{One decision for you}

The two clock designs are genuinely different bets:

\begin{itemize}
\item The \textbf{cascade} (stacked gears) is compact, digital, and closest to the
original ``nested loops at different scales'' image --- and Section~\ref{sec:struct}
shows it can be made robust.
\item The \textbf{incommensurate bank} (many unrelated speeds) is sturdier out of the
box but bulkier and shorter-range.
\end{itemize}

Which one to lean into shapes the next stage. The mild preference here is the cascade,
now that local coupling has made it robust --- it is the more elegant and more
faithful to the starting idea --- but it is your call.

\end{document}
